% !PRINTED 15
\weberpage{15}

The construction of a cut in the sequence of numbers always supplies
the proof for the existence of a number satisfying prescribed
requirements. Thus one obtains a cut $(A,B)$ by putting into $A$ all
and only those numbers whose square is smaller than a prescribed number
$\alpha$; to this cut there corresponds a definite number whose square
is equal to $\alpha$ and which is denoted by $\sqrt\alpha$. In this way
the existence of square roots is proved.

The numerical sequences introduced by G.\ \gesperrt{Cantor} for the
definition of irrational numbers can also be referred back to cuts.
\setcounter{footnote}{0}\footnote{Cantor, \emph{Ueber die Ausdehnung eines Satzes aus der
Theorie der trigonometrischen Reihen}. Mathematische Annalen, vol.~5
(1872); compare also Heine, \emph{Elemente der Functionenlehre}.
Journal für Mathematik, vol.~74 (1872).}

According to \gesperrt{Cantor}, by a numerical sequence one is to
understand any unbounded system of numbers ordered in a definite way:
\[
 S=x_1,x_2,x_3,x_4,\ldots
\]
having the property that there is a definite number $g$ below which
none of the numbers of $S$ sinks, and that, if $\delta$ is any chosen
number and $x_n,x_m$ are different elements of $S$, then $x_n-x_m$ or
$x_m-x_n$ always remains smaller than $\delta$ when $m$ and $n$ have
exceeded a sufficiently large number.

There are always numbers which are not exceeded by the numbers of such
a sequence; for, once $\delta$ has been chosen and $n$ suitably
determined, $x_m$, whatever the value of $m$, never exceeds the
greatest of the numbers $x_1,x_2,\ldots,x_n,x_n+\delta$. One now
obtains a cut $(A,B)$ by throwing into $B$ those numbers which, when
$n$ has a sufficiently high value, are no longer exceeded by any
$x_n$, and all other numbers, which are therefore exceeded by
infinitely many $x_n$, into $A$. If this cut is produced by the number
$\alpha$, then, however small $\varepsilon$ may be, there are always
infinitely many numbers $x_n$ between $\alpha-\varepsilon$ and
$\alpha$; and one may say that this number $\alpha$, completely
determined by $S$, is produced by the numerical sequence $S$. According
to \gesperrt{Cantor}, the numerical sequence $S$ is directly the
definition of the number $\alpha$. Of course one and the same
%CONT>>
